% -----------------------------------------------------------
\section{End (noun), Endless}
% -----------------------------------------------------------
	\label{END}
	
% % ----------------------
% \subsection{See also}
% % ----------------------

% ----------------------
\subsection{1828 American Dictionary of the English Language} % *1828
% ----------------------
	\label{END-1828}

	\begin{compactitem}
		\item[1] The extreme point of a line, or of anything that has more length than breadth; as the end of a house; the end of a table; the end of a finger; the end of a chain or rope. When bodies or figures have equal dimensions, or equal length and breadth, the extremities are called sides.
		\item[2] The extremity or last part, in general; the close or conclusion, applied to time.
		\item[3] The conclusion or cessation of an action.
		\item[4] The close or conclusion; as the end of a chapter.
		\item[5] Ultimate state or condition; final doom.
		\item[6] The point beyond which no progression can be made.
		\item[7] Final determination; conclusion of debate or deliberation.
		\item[8] Close of life; death; decease.
		\item[9] Cessation; period; close of a particular state of things; as the end of the world.
		\item[10] Limit; termination.
		\item[11] Destruction.
		\item[12] Cause of death; a destroyer.
		\item[13] Consequence; issue; result; conclusive event; conclusion.
		\item[14] A fragment or broken piece.
		\item[15] The ultimate point or thing at which one aims or directs his views; the object intended to be reached or accomplished by any action or scheme; purpose intended; scope; aim; drift; as private ends; public ends.
		\item[16] An end for on end upright; erect; as, his hair stands an end.
		\item[17] The ends of the earth, in scripture, are the remotest parts of the earth, or the inhabitants of those parts.
	\end{compactitem}

% % ----------------------
% \subsection{Oxford English Dictionary} % *OED
% % ----------------------
	% \label{END-OED}

	% \begin{compactitem}
		% \item[]
	% \end{compactitem}

% ----------------------
\subsection{Scriptural Usage} % *SCRIPTURES
% ----------------------
	\label{END-SCRIPTURES}
	
	Generally, I am not interested in any uses of this term that are similar to ``made an end,'' or that obviously refer to the extreme part of a physical object, like ``end of the rod of iron.''

	% % -----
	% \subsubsection{Old Testament} % *OT
	% % -----
		% \label{END-OT}
	
		% % --
		% \paragraph{KJV}
		% % --
			% \label{END-OT-KJV}
	
			% \begin{tabular}{ll}
				% Total occurrences in the OT & 
			% \end{tabular}
			
			% \begin{compactdesc}
				% \item[]
			% \end{compactdesc}
			
		% % --
		% \paragraph{Hebrew Bible}
		% % --
			% \label{END-HEBREW}
		
			% %
			% \subparagraph{\texorpdfstring{\he{sample word}}{}}
			% %
				% \label{PAGA} % whatever is in step bible without periods
			
				% \begin{compactdesc}
					% \item[HALOT , PDF ] \textbf{}
						% \begin{compactitem}
							% \item[1]
						% \end{compactitem}
					% \item[BDB , PDF ] \textbf{}
						% \begin{compactitem}
							% \item[1]
						% \end{compactitem}
					% \item[DCH PDF ] \textbf{}
						% \begin{compactitem}
							% \item[1]
						% \end{compactitem}
				% \end{compactdesc}
				
				% \begin{compactdesc}
					% \item[Some OT uses] \textbf{}
						% \begin{compactdesc}
							% \item[]
						% \end{compactdesc}
				% \end{compactdesc}
			
		% % --
		% \paragraph{Septuagint}
		% % --
			% \label{END-LXX}
		
			% %
			% \subparagraph{\texorpdfstring{\gr{sample word}}{}}
			% %
		
	% -----
	\subsubsection{New Testament} % *NT
	% -----
		\label{END-NT}
	
		% --
		\paragraph{KJV}
		% --
			\label{END-NT-KJV}
			
	
			\begin{tabular}{ll}
				Total occurrences in the NT & 60
			\end{tabular}
			
			\begin{compactdesc}
				\item[Matthew 10:22] And ye shall be hated of all men for my name's sake: but he that endureth to the end shall be saved.
				\item[Matthew 13:39] The enemy that sowed them is the devil; the harvest is the end of the world; and the reapers are the angels.
				\item[Matthew 13:40] As therefore the tares are gathered and burned in the fire; so shall it be in the end of this world.
				\item[Matthew 28:20] Teaching them to observe all things whatsoever I have commanded you: and, lo, I am with you alway, even unto the end of the world. Amen.
				\item[Mark 3:26] And if Satan rise up against himself, and be divided, he cannot stand, but hath an end.
				\item[Luke 1:33] And he shall reign over the house of Jacob for ever; and of his kingdom there shall be no end.
				\item[John 13:1] NOW before the feast of the passover, when Jesus knew that his hour was come that he should depart out of this world unto the Father, having loved his own which were in the world, he loved them unto the end.
				\item[John 18:37] Pilate therefore said unto him, Art thou a king then? Jesus answered, Thou sayest that I am a king. To this end was I born, and for this cause came I into the world, that I should bear witness unto the truth. Every one that is of the truth heareth my voice. (Here, ``end'' and ``cause'' are just \gr{τοῦτο}.)
				\item[Romans 4:16] Therefore it is of faith, that it might be by grace; to the end the promise might be sure to all the seed; not to that only which is of the law, but to that also which is of the faith of Abraham; who is the father of us all, (Here, ``to the end'' is a purpose clause, \gr{εἰς τὸ εἶναι}.)
				\item[Romans 6:21] What fruit had ye then in those things whereof ye are now ashamed? for the end of those things is death.
				\item[Romans 6:22] But now being made free from sin, and become servants to God, ye have your fruit unto holiness, and the end everlasting life.
				\item[Romans 10:4] For Christ is the end of the law for righteousness to every one that believeth.
				\item[Romans 14:9] For to this end Christ both died, and rose, and revived, that he might be Lord both of the dead and living. (``End'' here is \gr{τοῦτο} again, like John 18:37.)
				\item[1 Corinthians 15:24] Then cometh the end, when he shall have delivered up the kingdom to God, even the Father; when he shall have put down all rule and all authority and power.
				\item[2 Corinthians 3:13] And not as Moses, which put a vail over his face, that the children of Israel could not stedfastly look to the end of that which is abolished:
				\item[Philippians 3:18--19] (For many walk, of whom I have told you often, and now tell you even weeping, that they are the enemies of the cross of Christ: 19 Whose end is destruction, whose God is their belly, and whose glory is in their shame, who mind earthly things.)
				\item[1 Timothy 1:5] Now the end of the commandment is charity out of a pure heart, and of a good conscience, and of faith unfeigned:
				\item[Hebrews 7:3] Without father, without mother, without descent, having neither beginning of days, nor end of life; but made like unto the Son of God; abideth a priest continually.
				\item[James 5:11] Behold, we count them happy which endure. Ye have heard of the patience of Job, and have seen the end of the Lord; that the Lord is very pitiful, and of tender mercy.
				\item[1 Peter 1:9] Receiving the end of your faith, even the salvation of your souls.
				\item[2 Peter 2:20] For if after they have escaped the pollutions of the world through the knowledge of the Lord and Saviour Jesus Christ, they are again entangled therein, and overcome, the latter end is worse with them than the beginning.
				\item[Revelation 21:6] And he said unto me, It is done. I am Alpha and Omega, the beginning and the end. I will give unto him that is athirst of the fountain of the water of life freely.
			\end{compactdesc}
			
		% --
		\paragraph{Greek New Testament}
		% --
			\label{END-GREEK}
			
			%
			\subparagraph{\texorpdfstring{\gr{ἔσχατος}}{}}
			%
				\label{ESCHATOS}
			
				In the context of concept ``end,'' in 2 Peter 2:20 \gr{τὰ ἔσχατα} is just something like ``the things at the end,'' or ``the last things,'' as compared to the first ones. This would be meaning \#1 or \#2 in BDAG, I think.
			
				\begin{compactdesc}
					\item[BDAG 350b] \textbf{}
						\begin{compactitem}
							\item[1] pertaining to being at the farthest boundary of an area, \textit{farthest, last}
							\item[2] pertaining to being the final item in a series, \textit{least, last}
							\item[3] pertaining to furthest extremity in rank, value, or situation, \textit{last}
						\end{compactitem}
					\item[LSJ 699b, PDF 770] \textbf{}
						\begin{compactitem}
							\item[I.1] of space, \textit{farthest, uttermost, extreme}
							\item[I.2] of degree, \textit{uttermost, highest}
							\item[I.3] of persons, \textit{lowest, meanest}
							\item[I.4] of time, \textit{last}
						\end{compactitem}
				\end{compactdesc}
				
				\begin{compactdesc}
					\item[Some NT uses] \textbf{}
						\begin{compactdesc}
							\item[2 Peter 2:20] \gr{εἰ γὰρ ἀποφυγόντες τὰ μιάσματα τοῦ κόσμου ἐν ἐπιγνώσει τοῦ κυρίου καὶ σωτῆρος Ἰησοῦ Χριστοῦ τούτοις δὲ πάλιν ἐμπλακέντες ἡττῶνται, γέγονεν αὐτοῖς τὰ ἔσχατα χείρονα τῶν πρώτων.}
						\end{compactdesc}
				\end{compactdesc}
			
			%
			\subparagraph{\texorpdfstring{\gr{συντέλεια}}{}}
			%
				\label{SUNTELEIA}
				
				Notice how different the meanings are between BDAG and LSJ; it looks like only meaning III.1 in LSJ is related to the BDAG meaning. This word was commonly used for several things long before the NT but then its use and meaning narrowed considerably after that.
			
				\begin{compactdesc}
					\item[BDAG 866a] \textbf{}
						\begin{compactitem}
							\item a point of time marking completion of a duration, \textit{completion, close, end}
						\end{compactitem}
					\item[LSJ 1725b, PDF 1796] \textbf{}
						\begin{compactitem}
							\item[I.1] joint contribution for the public burdens
							\item[I.2] instruction
							\item[I.3] (compulsory) provision of recruits
							\item[II.1] a body of citizens who contributed jointly to bear public burdens
							\item[II.2] company
							\item[II.3] union of communities
							\item[III.1] the consummation of a scheme
							\item[IV] unjust gain
							\item[V] in grammar, \textit{completed action}
							\item[VI] reality (LSJ only cites a single example)
						\end{compactitem}
				\end{compactdesc}
				
				\begin{compactdesc}
					\item[Some NT uses] \textbf{}
						\begin{compactdesc}
							\item[Matthew 13:39] \gr{ὁ δὲ ἐχθρὸς ὁ σπείρας αὐτά ἐστιν ὁ διάβολος· ὁ δὲ θερισμὸς συντέλεια αἰῶνός ἐστιν, οἱ δὲ θερισταὶ ἄγγελοί εἰσιν.}
							\item[Matthew 13:40] \gr{ὥσπερ οὖν συλλέγεται τὰ ζιζάνια καὶ πυρὶ καίεται, οὕτως ἔσται ἐν τῇ συντελείᾳ τοῦ αἰῶνος·}
							\item[Matthew 28:20] \gr{διδάσκοντες αὐτοὺς τηρεῖν πάντα ὅσα ἐνετειλάμην ὑμῖν· καὶ ἰδοὺ ἐγὼ μεθ᾽ ὑμῶν εἰμι πάσας τὰς ἡμέρας ἕως τῆς συντελείας τοῦ αἰῶνος.}
						\end{compactdesc}
				\end{compactdesc}
		
			%
			\subparagraph{\texorpdfstring{\gr{τέλος}}{}}
			%
				\label{TELOS}
				
				The uses in Romans seem to clearly be meaning \#3 from BDAG; Philippians 3:18--19 as well; 1 Timothy 1:5 could be, though it could also be meaning 
			
				\begin{compactdesc}
					\item[BDAG 887a] \textbf{}
						\begin{compactitem}
							\item[1] a point of time marking the end of a duration, \textit{end, termination, cessation}
							\item[2] the last part of a process, \textit{close, conclusion}
							\item[3] the goal toward which a movement is being directed, \textit{end, goal, outcome}
							\item[4] last in a series, \textit{rest, remainder}
							\item[5] revenue obligation, \textit{(indirect) tax, toll-tax, customs duties}
						\end{compactitem}
					\item[LSJ 1772b, PDF 1843] \textbf{}
						\begin{compactitem}
							\item[I.1] coming to pass, performance, consummation
							\item[I.2] power of deciding, supreme power
							\item[I.3] magistracy, office
							\item[I.4] decision, doom
							\item[I.5] something done or ordered to be done, task, service, duty
							\item[I.6] plural, \textit{services} or \textit{offerings due} to the gods
							\item[I.7] service rendered by a citizen in the Solonian constitution
							\item[I.8] dues exacted by the state
							\item[I.9] financial means, expenditure
							\item[I.10] a military station or post with defined duties
							\item[I.11] troops or columns thereof
							\item[I.12] a territorial division
							\item[II.1.a] degree of completion or attainment
							\item[II.1.b] a length of time (or space), term, course
							\item[II.2] state of completion or maturity
							\item[II.3] to have reached the end of life, be dead
							\item[II.4] end, cessation
							\item[II.5] end of a word
							\item[III.1] achievement, attainment
							\item[III.2] winning-post, goal in a race
							\item[III.3.a] in philosophy, \textit{full realization, highest point, ideal}
							\item[III.3.b] the end or purpose of action
						\end{compactitem}
				\end{compactdesc}
				
				\begin{compactdesc}
					\item[Some NT uses] \textbf{}
						\begin{compactdesc}
							\item[Matthew 10:22] \gr{καὶ ἔσεσθε μισούμενοι ὑπὸ πάντων διὰ τὸ ὄνομά μου· ὁ δὲ ὑπομείνας εἰς τέλος οὗτος σωθήσεται.} (BDAG \#1, \#2 or \#3)
							\item[Mark 3:26] \gr{καὶ εἰ ὁ Σατανᾶς ἀνέστη ἐφ᾽ ἑαυτὸν καὶ ἐμερίσθη, οὐ δύναται στῆναι ἀλλὰ τέλος ἔχει.} (BDAG \#1)
							\item[Luke 1:33] \gr{καὶ βασιλεύσει ἐπὶ τὸν οἶκον Ἰακὼβ εἰς τοὺς αἰῶνας, καὶ τῆς βασιλείας αὐτοῦ οὐκ ἔσται τέλος.} (BDAG \#1 or \#2)
							\item[John 13:1] \gr{Πρὸ δὲ τῆς ἑορτῆς τοῦ πάσχα εἰδὼς ὁ Ἰησοῦς ὅτι ἦλθεν αὐτοῦ ἡ ὥρα ἵνα μεταβῇ ἐκ τοῦ κόσμου τούτου πρὸς τὸν πατέρα ἀγαπήσας τοὺς ἰδίους τοὺς ἐν τῷ κόσμῳ εἰς τέλος ἠγάπησεν αὐτούς.} (BDAG \#1 or \#2)
							\item[Romans 6:21] \gr{τίνα οὖν καρπὸν εἴχετε τότε ἐφ᾽ οἷς νῦν ἐπαισχύνεσθε; τὸ γὰρ τέλος ἐκείνων θάνατος·} (BDAG \#3)
							\item[Romans 6:22] \gr{νυνὶ δέ, ἐλευθερωθέντες ἀπὸ τῆς ἁμαρτίας δουλωθέντες δὲ τῷ θεῷ, ἔχετε τὸν καρπὸν ὑμῶν εἰς ἁγιασμόν, τὸ δὲ τέλος ζωὴν αἰώνιον.} (BDAG \#3)
							\item[Romans 10:4] \gr{τέλος γὰρ νόμου Χριστὸς εἰς δικαιοσύνην παντὶ τῷ πιστεύοντι.} (BDAG \#3, possibly \#2)
							\item[1 Corinthians 15:24] \gr{εἶτα τὸ τέλος, ὅταν παραδιδῷ τὴν βασιλείαν τῷ θεῷ καὶ πατρί, ὅταν καταργήσῃ πᾶσαν ἀρχὴν καὶ πᾶσαν ἐξουσίαν καὶ δύναμιν,} (BDAG \#2)
							\item[2 Corinthians 3:13] \gr{καὶ οὐ καθάπερ Μωϋσῆς ἐτίθει κάλυμμα ἐπὶ τὸ πρόσωπον αὐτοῦ, πρὸς τὸ μὴ ἀτενίσαι τοὺς υἱοὺς Ἰσραὴλ εἰς τὸ τέλος τοῦ καταργουμένου.} (BDAG \#1, \#2, or \#3)
							\item[Philippians 3:18--19] \gr{πολλοὶ γὰρ περιπατοῦσιν οὓς πολλάκις ἔλεγον ὑμῖν, νῦν δὲ καὶ κλαίων λέγω, τοὺς ἐχθροὺς τοῦ σταυροῦ τοῦ Χριστοῦ,} 19 \gr{ὧν τὸ τέλος ἀπώλεια, ὧν ὁ θεὸς ἡ κοιλία καὶ ἡ δόξα ἐν τῇ αἰσχύνῃ αὐτῶν, οἱ τὰ ἐπίγεια φρονοῦντες.} (BDAG \#2 or \#3)
							\item[1 Timothy 1:5] \gr{τὸ δὲ τέλος τῆς παραγγελίας ἐστὶν ἀγάπη ἐκ καθαρᾶς καρδίας καὶ συνειδήσεως ἀγαθῆς καὶ πίστεως ἀνυποκρίτου,} (BDAG \#3)
							\item[Hebrews 7:3] \gr{ἀπάτωρ, ἀμήτωρ, ἀγενεαλόγητος, μήτε ἀρχὴν ἡμερῶν μήτε ζωῆς τέλος ἔχων, ἀφωμοιωμένος δὲ τῷ υἱῷ τοῦ θεοῦ, μένει ἱερεὺς εἰς τὸ διηνεκές.} (BDAG \#1 or \#2)
							\item[James 5:11] \gr{ἰδοὺ μακαρίζομεν τοὺς ὑπομείναντας· τὴν ὑπομονὴν Ἰὼβ ἠκούσατε, καὶ τὸ τέλος κυρίου εἴδετε, ὅτι πολύσπλαγχνός ἐστιν ὁ κύριος καὶ οἰκτίρμων.} (BDAG \#3)
							\item[1 Peter 1:9] \gr{κομιζόμενοι τὸ τέλος τῆς πίστεως ὑμῶν σωτηρίαν ψυχῶν.} (BDAG \#3)
							\item[Revelation 21:6] \gr{καὶ εἶπέν μοι· Γέγοναν. ἐγὼ τὸ Ἄλφα καὶ τὸ Ὦ, ἡ ἀρχὴ καὶ τὸ τέλος. ἐγὼ τῷ διψῶντι δώσω ἐκ τῆς πηγῆς τοῦ ὕδατος τῆς ζωῆς δωρεάν.} (BDAG \#1, \#2, or \#3)
						\end{compactdesc}
				\end{compactdesc}
		
	% -----
	\subsubsection{Book of Mormon} % *BOM
	% -----
		\label{END-BOM}
	
		\begin{tabular}{ll}
			Total occurrences in the BOM & 191
		\end{tabular}
		
		\begin{compactdesc}
			\item[1 Nephi 13:37] And blessed are they who shall seek to bring forth my Zion at that day, for they shall have the gift and the power of the Holy Ghost; and if they endure unto the end they shall be lifted up at the last day, and shall be saved in the everlasting kingdom of the Lamb; and whoso shall publish peace, yea, tidings of great joy, how beautiful upon the mountains shall they be.
				\textit{(See the \hyperref[END-PHRASES]{key phrases section} below for commentary on the phrase ``endure to the end.'')}
			\item[1 Nephi 14:3] And that great pit, which hath been digged for them by that great and abominable church, which was founded by the devil and his children, that he might lead away the souls of men down to hell---yea, that great pit which hath been digged for the destruction of men shall be filled by those who digged it, unto their utter destruction, saith the Lamb of God; not the destruction of the soul, save it be the casting of it into that hell which hath no end.
			\item[1 Nephi 15:30] And I said unto them that our father also saw that the justice of God did also divide the wicked from the righteous; and the brightness thereof was like unto the brightness of a flaming fire, which ascendeth up unto God forever and ever, and hath no end.
			\item[2 Nephi 2:7] Behold, he offereth himself a sacrifice for sin, to answer the ends of the law, unto all those who have a broken heart and a contrite spirit; and unto none else can the ends of the law be answered.
				\textit{(See the \hyperref[ANSWER-1828]{definitions of ``answer''} in the 1828 dictionary; I think the meaning here has to be definition \#2, \#3, or \#6 from that list. In fact I think it has to be \#3 or \#6. See \hyperref[2NEPHI-2-7]{2 Nephi 2:7} for more commentary. See also the \hyperref[END-PHRASES]{key phrases section} for ``ends of the law''; the key verse in the NT is \hyperref[ROMANS-10-4]{Romans 10:4}, and the phrase is \gr{τέλος γὰρ νόμου Χριστὸς}, though I'm still not sure how much help that phrase is here.)}
			\item[2 Nephi 2:10] And because of the intercession for all, all men come unto God; wherefore, they stand in the presence of him, to be judged of him according to the truth and holiness which is in him. Wherefore, the ends of the law which the Holy One hath given, unto the inflicting of the punishment which is affixed, which punishment that is affixed is in opposition to that of the happiness which is affixed, to answer the ends of the atonement--- \textit{(I think the way to understand ``ends'' here is as ``purposes.'' The ``ends'' of the law are the possible outcomes---and so in that sense maybe there is a bit of \gr{συντέλεια} in there as well---and the atonement, as we read at the end, also has these ``ends,'' or purposes, or things that are brought about by the atonement. And in accordance with the first use, I think that the possible ends or outcomes are happiness and misery. But I also think it's important to remember that the purpose or end of the atonement is to \textbf{allow us to experience the things which our actions have already said we wish to recieve}. So I wouldn't want to look at it is the atonement inflicting a punishment of misery or compelling us to be happy; instead, I think the atonement makes it possible to get a fuller version of what we really truly want or are comfortable with, in the spirit of \hyperref[DC88-32]{DC 88:32}.)}
			\item[2 Nephi 2:12] Wherefore, it must needs have been created for a thing of naught; wherefore there would have been no purpose in the end of its creation. Wherefore, this thing must needs destroy the wisdom of God and his eternal purposes, and also the power, and the mercy, and the justice of God. \textit{(So how do we read this verse, given what I said about 2 Nephi 2:10 above? I think we can read it in a similar way. Our creation was an end in itself---it was a good or purpose that was worth achieving by itself. But it wasn't the final end---there was further purpose as well, which is for us to eventually be exalted. There are ``eternal purposes in the end of man,'' as \hyperref[2NEPHI-2-15]{2 Nephi 2:15} says below.)}
			\item[2 Nephi 2:15] And to bring about his eternal purposes in the end of man, after he had created our first parents, and the beasts of the field and the fowls of the air, and in fine, all things which are created, it must needs be that there was an opposition; even the forbidden fruit in opposition to the tree of life; the one being sweet and the other bitter. \textit{(See \hyperref[2NEPHI-2-12]{2 Nephi 2:12} above; there are eternal purposes in the end of man, or purposes that go beyond the mere creation of human beings. There is more to happen and be done.)}
			\item[2 Nephi 2:22] And now, behold, if Adam had not transgressed he would not have fallen, but he would have remained in the garden of Eden. And all things which were created must have remained in the same state in which they were after they were created; and they must have remained forever, and had no end. \textit{(Here is a case in 2 Nephi where ``end'' does seem to be more of a straight \gr{συντέλεια}. Especially given that it follows the ``remained forever'' bit immediately, there would be no end to their duration. But notice how important this is in light of ends as purposes. If ``end'' here is just more of \gr{συντέλεια}, then there \textbf{must} be an end (\gr{συντέλεια}) in order for the larger purposes to be accomplished.)}
			\item[2 Nephi 9:16] And assuredly, as the Lord liveth, for the Lord God hath spoken it, and it is his eternal word, which cannot pass away, that they who are righteous shall be righteous still, and they who are filthy shall be filthy still; wherefore, they who are filthy are the devil and his angels; and they shall go away into everlasting fire, prepared for them; and their torment is as a lake of fire and brimstone, whose flame ascendeth up forever and ever and has no end.
			\item[2 Nephi 11:4] Behold, my soul delighteth in proving unto my people the truth of the coming of Christ; for, for this end hath the law of Moses been given; and all things which have been given of God from the beginning of the world, unto man, are the typifying of him.
			\item[2 Nephi 25:25] For, for this end was the law given; wherefore the law hath become dead unto us, and we are made alive in Christ because of our faith; yet we keep the law because of the commandments.
			\item[2 Nephi 25:27] Wherefore, we speak concerning the law that our children may know the deadness of the law; and they, by knowing the deadness of the law, may look forward unto that life which is in Christ, and know for what end the law was given. And after the law is fulfilled in Christ, that they need not harden their hearts against him when the law ought to be done away.
			\item[2 Nephi 29:9] And I do this that I may prove unto many that I am the same yesterday, today, and forever; and that I speak forth my words according to mine own pleasure. And because that I have spoken one word ye need not suppose that I cannot speak another; for my work is not yet finished; neither shall it be until the end of man, neither from that time henceforth and forever. \textit{(A powerful verse for ``end.'' God's work will not be finished until ``the end of man''---could be taken in the purpose sense but also in the temporal, \gr{συντέλεια} sense. If you take it in the former sense, then it would be something like, ``until the purpose of/in man has been accomplished.'' But if you take it in the latter temporal sense, then there is a time when man will end, or be no more, since we will have been transformed into some other kind of being; even then, God's work will not be finished. A powerful verse for thinking about exaltation.)}
			\item[2 Nephi 31:20] Wherefore, ye must press forward with a steadfastness in Christ, having a perfect brightness of hope, and a love of God and of all men. Wherefore, if ye shall press forward, feasting upon the word of Christ, and endure to the end, behold, thus saith the Father: Ye shall have eternal life.
			\item[2 Nephi 31:21] And now, behold, my beloved brethren, this is the way; and there is none other way nor name given under heaven whereby man can be saved in the kingdom of God. And now, behold, this is the doctrine of Christ, and the only and true doctrine of the Father, and of the Son, and of the Holy Ghost, which is one God, without end. Amen.
			\item[2 Nephi 33:9] I also have charity for the Gentiles. But behold, for none of these can I hope except they shall be reconciled unto Christ, and enter into the narrow gate, and walk in the strait path which leads to life, and continue in the path until the end of the day of probation.
			\item[Jacob 2:21] Do ye not suppose that such things are abominable unto him who created all flesh? And the one being is as precious in his sight as the other. And all flesh is of the dust; and for the selfsame end hath he created them, that they should keep his commandments and glorify him forever.
			\item[Omni 1:26] And now, my beloved brethren, I would that ye should come unto Christ, who is the Holy One of Israel, and partake of his salvation, and the power of his redemption.  Yea, come unto him, and offer your whole souls as an offering unto him, and continue in fasting and praying, and endure to the end; and as the Lord liveth ye will be saved.
			\item[Mosiah 4:6] I say unto you, if ye have come to a knowledge of the goodness of God, and his matchless power, and his wisdom, and his patience, and his long-suffering towards the children of men; and also, the atonement which has been prepared from the foundation of the world, that thereby salvation might come to him that should put his trust in the Lord, and should be diligent in keeping his commandments, and continue in the faith even unto the end of his life, I mean the life of the mortal body--- \textit{(See also Mosiah 4:30, 5:8, and Alma 34:33.)}
			\item[Alma 12:27] But behold, it was not so; but it was appointed unto men that they must die; and after death, they must come to judgment, even that same judgment of which we have spoken, which is the end. \textit{(Another key verse---the judgment after death is ``the end.'' The end of the probationary period, maybe? Though it would be important to connect this with what Alma says later about the judgment and the spirit world in Alma 40; see also Alma 34:33 below.)}
			\item[Alma 34:33] And now, as I said unto you before, as ye have had so many witnesses, therefore, I beseech of you that ye do not procrastinate the day of your repentance until the end; for after this day of life, which is given us to prepare for eternity, behold, if we do not improve our time while in this life, then cometh the night of darkness wherein there can be no labor performed. \textit{(This verse is interesting in light of Alma 12:27, where ``the end'' is the judgment after death. Putting them together, the instruction is to not procrastinate the day of repentance until that day of judgment. So the end of ``this day of life''---and that phrasing helps you combine these two verses with Mosiah 4:6---is the judgment, according to what we read here.)}
			\item[Mormon 7:7]  And he hath brought to pass the redemption of the world, whereby he that is found guiltless before him at the judgment day hath it given unto him to dwell in the presence of God in his kingdom, to sing ceaseless praises with the choirs above, unto the Father, and unto the Son, and unto the Holy Ghost, which are one God, in a state of happiness which hath no end.
			\item[Ether 13:8] Wherefore, the remnant of the house of Joseph shall be built upon this land; and it shall be a land of their inheritance; and they shall build up a holy city unto the Lord, like unto the Jerusalem of old; and they shall no more be confounded, until the end come when the earth shall pass away.
			\item[Moroni 7:28] For he hath answered the ends of the law, and he claimeth all those who have faith in him; and they who have faith in him will cleave unto every good thing; wherefore he advocateth the cause of the children of men; and he dwelleth eternally in the heavens.
			\item[Moroni 8:26] And the remission of sins bringeth meekness, and lowliness of heart; and because of meekness and lowliness of heart cometh the visitation of the Holy Ghost, which Comforter filleth with hope and perfect love, which love endureth by diligence unto prayer, until the end shall come, when all the saints shall dwell with God.
		\end{compactdesc}
		
	% -----
	\subsubsection{Doctrine and Covenants} % *DC
	% -----
		\label{END-DC}
	
		\begin{tabular}{ll}
			Total occurrences in the DC & 91
		\end{tabular}
		
		\begin{compactdesc}
			\item[DC 19:6] Nevertheless, it is not written that there shall be no end to this torment, but it is written \textit{endless torment}.
		\end{compactdesc}
		
	% -----
	\subsubsection{Pearl of Great Price} % *PGP
	% -----
		\label{END-PGP}
	
		\begin{tabular}{ll}
			Total occurrences in the PGP & 26
		\end{tabular}
		
		\begin{compactdesc}
			\item[]
		\end{compactdesc}
		
% ----------------------
\subsection{Key Phrases} % *KEYPHRASES *PHRASES
% ----------------------
	\label{END-PHRASES}

	\begin{compactdesc}
	
		\item[beginning and the end] Revelation 1:8, 21:6, 22:13; 2 Nephi 27:7; Alma 11:39; 3 Nephi 9:18; DC 19:1; 35:1; 38:1; 45:7; 49:12; 54:1; 61:1; 84:120; 95:7; Moses 2:1
			\begin{compactdesc}
				\item[Bible anchor verses] Revelation 1:8, 21:6, 22:13
					\begin{compactdesc}
						\item[Revelation 1:8] --- (the critical text removes ``the beginning and the ending'' here, with a certainty rating of ``A.'')
						\item[Revelation 21:6] \gr{ἡ ἀρχὴ καὶ τὸ τέλος}
						\item[Revelation 22:13] \gr{ἡ ἀρχὴ καὶ τὸ τέλος}
					\end{compactdesc}
				\item[Commentary] For the verses in Revelation 21 and 22, probably meanings \#1--\#3 of BDAG are in play, which makes them very rich. So I think we should see the uses in the BOM and DC in the same way. Those verses are speaking about Christ, so for the ``beginning'' part, you might think of Christ as the \hyperref[FIRSTFRUITS]{firstfruits}, and for the end, you might think of how the culmination of our process in the plan of salvation is to become like Christ.
			\end{compactdesc}
			
		\item[beginning of days, end of years] \textbf{5} | \hyperref[ALMA-13-7]{Alma 13:7}, \hyperref[ALMA-13-9]{9}; \hyperref[DC84-17]{DC 84:17}; \hyperref[MOSES-1-3]{Moses 1:3}; \hyperref[MOSES-6-67]{Moses 6:67}
			\begin{compactdesc}
				\item[Bible anchor verses] It isn't exact, but see \hyperref[HEBREWS-7-3]{Hebrews 7:3}
				\item[Commentary] Why ``days'' and ``years,'' two different measurements of time? There could be a lot of explanations, but I think of how \textit{days} are the measure of time in the creation story, and we use \textit{years} to mark how old we are. Days and years are also related in \hyperref[GENESIS-1-14]{Genesis 1:14}, \hyperref[JOB-10-5]{Job 10:5}, \hyperref[PSALMS-77-5]{Psalms 77:5}, \hyperref[PSALMS-90-10]{Psalms 90:10}, \hyperref[2PETER-3-8]{2 Peter 3:8}, and \hyperref[ABRAHAM-3-5]{Abraham 3:5}, among other places.
			\end{compactdesc}
			
		\item[end from the beginning] Isaiah 46:10; Abraham 2:8
			\begin{compactdesc}
				\item[Bible anchor verses] Isaiah 46:10
					\begin{compactdesc}
						\item[Isaiah 46:10] \he{mere'+siyt 'a.ha:riyt} (See \hyperref[RESHIT]{\he{re'+siyt}} at \hyperref[FIRSTFRUITS]{firstfruits})
					\end{compactdesc}
			\end{compactdesc}
			
		\item[end(s) of the earth] Deuteronomy 13:7 (2), 28:49, 28:64, 33:17; 1 Samuel 2:10; Job 28:24, 37:3, 38:13; Psalms 46:9, 48:10, 59:13, 61:2, 65:5, 67:7, 72:8, 98:3, 135:7; Proverbs 17:24, 30:4; Isaiah 5:26, 26:15, 40:28, 41:5, 41:9, 42:10, 43:6, 45:22, 48:20, 49:6, 52:10; Jeremiah 10:13, 16:19, 25:31, 25:33 (2), 51:16; Daniel 4:22; Micah 5:4; Zechariah 9:10; Acts 13:47; 1 Nephi 20:20, 21:6; 2 Nephi 15:26, 24:2, 26:25, 27:11, 29:2, 33:10, 33:13; Mosiah 12:24, 15:31; Alma 5:50, 29:7; 3 Nephi 9:22, 11:41, 16:20, 20:35, 27:20; Mormon 3:18, 3:22, 9:21, 9:25; Ether 3:25, 4:18; Moroni 7:34, 10:24; DC 1:11; 38:5; 43:21, 43:31; 45:26, 45:49; 58:45; 65:1, 65:2; 72:21; 88:101; 90:9; 105:39; 107:42; 109:23, 109:57, 109:67; 112:4; 122:1; 133:3; JS-M 1:55
			\begin{compactdesc}
				\item[Bible anchor verses] Job 28:24; Psalms 72:8; Isaiah 40:28, 45:22; Zechariah 9:10; Acts 13:47
					\begin{compactdesc}
						\item[Job 28:24] \he{liq:.sOt--hA'ArE.S}
						\item[LXX Job 28:24] \gr{αὐτὸς γὰρ τὴν ὑπ' οὐρανὸν πᾶσαν ἐφορᾷ} (This is pretty different from the Hebrew, and doesn't have a ``ends of the earth''-type phrase.)
						\item[Psalms 72:8] \he{`ad--'ap:sey--'ArE.S}
						\item[LXX Psalms 72:8] \gr{ἕως περάτων τῆς οἰκουμένης}
						\item[Isaiah 40:28] \he{q:.sOt hA'ArE.S}
						\item[LXX Isaiah 40:28] \gr{τὰ ἄκρα τῆς γῆς}
						\item[Isaiah 45:22] \he{k*Al--'ap:sey--'ArE.S}
						\item[LXX Isaiah 45:22] \gr{οἱ ἀπ' ἐσχάτου τῆς γῆς}
						\item[Zechariah 9:10] \he{`ad--'ap:sey--'ArE.S}
						\item[LXX Isaiah 45:22] \gr{καὶ κατάρξει ὑδάτων ἕως θαλάσσης καὶ ποταμῶν διεκβολὰς γῆς}
						\item[Acts 13:47] \gr{ἕως ἐσχάτου τῆς γῆς}
					\end{compactdesc}
			\end{compactdesc}
			
		\item[end(s) of the law] Romans 10:4; 2 Nephi 2:7 (2), 2:10; Moroni 7:28
			\begin{compactdesc}
				\item[Bible anchor verses] Romans 10:4
					\begin{compactdesc}
						\item[Romans 10:4] \gr{τέλος γὰρ νόμου Χριστὸς}
					\end{compactdesc}
			\end{compactdesc}
			
		\item[end(s) of the world] Psalms 19:4, 22:27; Isaiah 62:11; Matthew 13:39, 13:49, 24:3, 28:20; Romans 10:18; 1 Corinthians 10:11; Hebrews 9:26; 1 Nephi 14:22; Mosiah 4:7; DC 1:23; 19:3; 45:22; 132:49; Moses 6:7, 7:67; JS-M 1:4 (2)
		
		\item[endless torment] \textbf{8} | \hyperref[2NEPHI-9-19]{2 Nephi 9:19}, \hyperref[2NEPHI-9-26]{26}; \hyperref[2NEPHI-28-23]{2 Nephi 28:23}; \hyperref[JACOB-6-10]{Jacob 6:10}; \hyperref[MOSIAH-3-25]{Mosiah 3:25}; \hyperref[MOSIAH-28-3]{Mosiah 28:3}; \hyperref[MORONI-8-21]{Moroni 8:21}; \hyperref[DC19-6]{DC 19:6}
			\begin{compactdesc}
				\item[Bible anchor verses] ---
				\item[Commentary] See notes at and around \hyperref[DC19-6]{DC 19:6}.
			\end{compactdesc}
		
		\item[endure to the end] Matthew 10:22; 24:13; Mark 13:13; 1 Nephi 13:37; 22:31; 2 Nephi 9:24; 31:15, 31:16, 31:20; 33:3; Omni 1:26; Alma 32:13, 32:15; 38:2; 3 Nephi 15:9 (2); 27:6, 27:16; Mormon 9:29; DC 14:7; 18:22
			\begin{compactdesc}
				\item[Bible anchor verses] Matthew 10:22
					\begin{compactdesc}
						\item[Matthew 10:22] \gr{ὁ δὲ ὑπομείνας εἰς τέλος}
					\end{compactdesc}
				\item[Commentary] In this phrase, ``endure'' is fairly clear, though it would help to be clearer about what it means. But ``end'' is the key word---what is the ``end'' being spoken of? Just how long do we have to endure? The Book of Mormon suggests that ``end'' here has the meaning of \gr{συντέλεια} and related meanings of \gr{τέλος}, rather than the purpose/Aristotelian meanings of \gr{τέλος}. The key verse is \hyperref[MOSIAH-4-6]{Mosiah 4:6}, which speaks of people who ``should be diligent in keeping his commandments, and continue in the faith even unto the end \textbf{of his life, I mean the life of the mortal body}.'' We're talking about continuing in the faith, which seems similar enough to ``enduring,'' and then we are told that the enduring needs to continue until the end of one's mortal life. This verse really couldn't be clearer. Mosiah 4:30 and 5:8 say similar things (``even unto the end of your lives'' and ``should be obedient unto the end of your lives,'' respectively). This interpretation of the phrase ``endure to the end'' receives some support from \hyperref[ALMA-34-32]{Alma 34:32--34}, which speaks multiple times of ``this life'' and then clarifies in verse 34 that it is talking about the life of the mortal body, since Amulek talks about a certain spirit possessing our \textit{bodies} at the time that we go out of this (mortal) life. Another possible source of support is \hyperref[ALMA-40-11]{Alma 40:11--12}, where we read that the spirits of righteous people enter a state of rest and peace after they die. It may be that once you get to this state, no more enduring is needed. So the meaning of the phrase ``endure to the end,'' at least in the Book of Mormon, appears to be enduring until the end of our mortal lives. I think that JS later adds on to this view with the idea that we can actually do salvation-worthy things after death, contra what \hyperref[ALMA-34-33]{Alma 34:33} suggests. I think BY adds on to this as well with the view that we are constantly changing and either progressing or regressing, and that these changes for good or bad happen at every stage of life, no matter what. I think this makes more sense in light of what else we know about the plan of salvation, and is more faithful to the subversive and universalist nature of the restoration, since the ``endure until the end of the race, then rest if you're righteous'' view feels like a very traditional one.
			\end{compactdesc}
			
		\item[world(s) without end] Isaiah 45:17; Ephesians 3:21; DC 76:112
			\begin{compactdesc}
				\item[Bible anchor verses] Isaiah 45:17; Ephesians 3:21
					\begin{compactdesc}
						\item[Isaiah 45:17] \he{`ad--`Ol:mey `ad}
						\item[LXX Isaiah 45:17] \gr{ἕως τοῦ αἰῶνος}
						\item[Ephesians 3:21] \gr{εἰς πάσας τὰς γενεὰς τοῦ αἰῶνος τῶν αἰώνων}
					\end{compactdesc}
			\end{compactdesc}
	\end{compactdesc}
	
% % ----------------------
% \subsection{Bible Dictionaries} % *DICTIONARIES
% % ----------------------
	% \label{END-BD}

% % ----------------------
% \subsection{Restoration Usage} % *RESTORATION
% % ----------------------
	% \label{END-RESTORATION}

	% % -----
	% \subsubsection{Joseph Smith} % JS
	% % -----
		% \label{END-JS}
	
		% \begin{compactdesc}
			% \item[DATE/SOURCE]
		% \end{compactdesc}
	
	% % -----
	% \subsubsection{Brigham Young} % BY
	% % -----
		% \label{END-BY}
	
		% \begin{compactdesc}
			% \item[DATE/SOURCE]
		% \end{compactdesc}
	
% % ----------------------
% \subsection{Commentary and Studies} % *COMMENTARY *STUDIES
% % ----------------------
	% \label{END-COMMENTARY}

	% % -----
	% \subsubsection{Key Points}
	% % -----
		% \label{END-KEYS}
	
		% \begin{compactitem}
			% \item
		% \end{compactitem}